\section{Показатели качества и~ресурсов}\label{sec:metrics}
\subsection{Успех всей процедуры и~её стоимость}
Пусть набор содержит $N$ задач, для каждой из которых назначено $R$ повторных генераций. Индексом $i$ обозначим задачу, а~индексом $r$ её повтор. В~выполненных сериях каждый повтор представлял собой новый запуск выбранного способа решения. Такой запуск включал один или несколько вызовов модели. Его номер служил для учёта запусков. Случайное состояние модели на стороне поставщика по этому номеру не задавалось, поэтому повторение запроса не гарантировало одинакового ответа.

Итоговую программу проверяет независимый оценщик. Его скрытые тесты не передаются модели. Результат проверки обозначим через $q_{ir}\in\{0,1\}$: единица соответствует успеху, ноль неуспеху. Этот результат может отличаться от результата проверки $\mathcal{V}_{z_i}$, доступной самой процедуре во время работы.

Определим долю успешных запусков $Q$, средний расход токенов $\overline{T}$ и~среднюю стоимость обращений к~моделям $\overline{C}$:
\begin{formulakey}
\begin{equation}
Q=\frac{1}{NR}\sum_{i=1}^{N}\sum_{r=1}^{R}q_{ir},\quad
\overline{T}=\frac{1}{NR}\sum_{i,r}T_{ir},\quad
\overline{C}=\frac{1}{NR}\sum_{i,r}C_{ir}.
\end{equation}
\begin{symbols}
\sym{Q}{доля запусков, в~которых программа прошла все итоговые тесты}
\sym{\overline T,\,\overline C}{средние расход токенов и~стоимость одного запуска}
\sym{N,\ R,\ NR}{число задач, число повторов каждой и~общее число запусков одного способа}
\sym{i,\ r}{номер задачи и~номер её повтора}
\sym{q_{ir}}{результат проверки программы на задаче~$i$ в~повторе~$r$, 1 при успехе и~0 при неуспехе}
\sym{T_{ir},\,C_{ir}}{все токены и~стоимость данного запуска, включая неудачные попытки}
\sym{\sum_{i,r}}{сумма по всем задачам и~всем их повторам}
\end{symbols}
\end{formulakey}
Если часть наблюдений отсутствует, вместе с~показателем сообщаем число доступных запусков и~число пропусков.

Показатели эффективности определим как
\begin{formulakey}
\begin{equation}
U_{\mathrm{tok}}=\frac{Q}{\overline{T}},\qquad
U_{\$}=\frac{Q}{\overline{C}}.
\label{eq:efficiency}
\end{equation}
\begin{symbols}
\sym{U_{\mathrm{tok}},\,U_{\$}}{число успешных запусков на токен и~на денежную единицу соответственно}
\sym{Q,\,\overline T,\,\overline C}{доля успехов, среднее число токенов и~средняя стоимость запуска}
\end{symbols}
\end{formulakey}
При положительных знаменателях $U_{\mathrm{tok}}$ показывает число успешных запусков на один затраченный токен, а~$U_{\$}$ на одну денежную единицу. Для расчёта используются суммарные успехи и~суммарные затраты всех запусков. Усреднение отдельных отношений дало бы другой показатель.

Вместе с~эффективностью приводятся доля успехов $Q$ и~расходы. Даже высокий показатель эффективности может сопровождаться неприемлемо низким качеством. Для наборов разных типов результаты сначала показываются раздельно. Веса для общего среднего выбираются до просмотра результатов.

Величина $Q$ относится к~результату всей процедуры, которая может включать несколько кандидатов или исправлений. Выражение <<успех в~пределах бюджета>> используется, когда бюджет явно задан и~контролируется. Оценка решения с~первой попытки требует отдельного расчёта. Показатель pass@1 относится к~одному сгенерированному кандидату без отбора по скрытым тестам~\cite{chen2021humaneval}.

\subsection{Дополнительные проверки и~свойства кода}
Дополнительную проверку программы можно разделить на тесты, статический анализ и~проверки безопасности. Их результаты обозначим через $t_{ir}$, $a_{ir}$ и~$u_{ir}$ соответственно. Каждый результат принимает значение $1$ при успешном прохождении проверки и~ноль при неуспехе. Доля запусков, прошедших все три вида проверки, равна
\begin{formulakey}
\begin{equation}
\mathrm{VCR}=\frac{1}{NR}\sum_{i,r}\I[t_{ir}=1\land a_{ir}=1\land u_{ir}=1].
\end{equation}
\begin{symbols}
\sym{\mathrm{VCR}}{доля запусков, в~которых программа прошла все три вида проверки}
\sym{N,\,R,\,NR}{число задач, число повторов на задачу и~общее число запусков одного способа}
\sym{i,\,r,\,\sum_{i,r}}{номер задачи, номер повтора и~сумма по всем таким парам}
\sym{t_{ir},\,a_{ir},\,u_{ir}}{результаты тестов, статического анализа и~проверок безопасности. В~каждом поле 1 означает успех, 0 означает неуспех}
\sym{\land,\,\I[\cdot]}{логическое <<и>> и~индикатор. Слагаемое равно 1, только когда все три результата равны 1}
\end{symbols}
\end{formulakey}
Для расчёта нужны заранее определённые проверки и~результаты каждого вида. Если часть проверок отсутствует, показатель по всей выборке не рассчитывается.

Для описания свойств программы введён дополнительный индекс $M$. Он объединяет четыре признака. Цикломатическая сложность характеризует ветвление программы: она равна числу независимых путей в~графе возможного выполнения команд~\cite{mccabe1976}. Остальные признаки описывают долю повторяющегося кода, плотность предупреждений анализатора и~долю документированных публичных функций.

Этот индекс характеризует измеряемые свойства текста программы. Работа Sepidband и~соавторов~\cite{sepidband2025complexity} исследует связь метрик сложности с~результатами генерации. Предлагаемые здесь веса признаков в~ней не проверялись.

Среднюю цикломатическую сложность функций одного кандидата обозначим через $\mathrm{CC}$. Долю дублирующихся строк обозначим через $\mathrm{Dup}$, число предупреждений на тысячу непустых строк через $\mathrm{Warn}$, а~долю документированных публичных функций через $\mathrm{Doc}$. Для приведения признаков к~общей шкале заранее задаются положительные масштабы $c_{\mathrm{cc}}$ и~$c_{\mathrm{warn}}$. Их выбирают на отдельном наборе. Нормированные значения вычисляются по формулам:
\begin{formulakey}
\begin{align}
\widehat{\mathrm{CC}}&=\min(1,\mathrm{CC}/c_{\mathrm{cc}}), &
\widehat{\mathrm{Warn}}&=\min(1,\mathrm{Warn}/c_{\mathrm{warn}}),\\
\widehat{\mathrm{Dup}}&=\mathrm{Dup}, &
\widehat{\mathrm{Doc}}&=\mathrm{Doc}.
\end{align}
\begin{symbols}
\sym{\mathrm{CC},\,\mathrm{Warn}}{средняя цикломатическая сложность и~число предупреждений на тысячу непустых строк}
\sym{c_{\mathrm{cc}},\,c_{\mathrm{warn}}}{положительные масштабы нормировки для сложности и~предупреждений}
\sym{\mathrm{Dup},\,\mathrm{Doc}}{доли повторяющихся строк и~документированных публичных функций}
\sym{\widehat{\mathrm{CC}},\,\widehat{\mathrm{Warn}},\,\widehat{\mathrm{Dup}},\,\widehat{\mathrm{Doc}}}{соответствующие признаки на шкале от 0 до~1. Операция $\min(1,x)$ ограничивает значение сверху единицей}
\end{symbols}
\end{formulakey}
Все четыре значения лежат в~$[0,1]$. Метод обнаружения дубликатов, список предупреждений и~критерий наличия документации фиксируются заранее. Если хотя бы один признак недоступен, например из-за отсутствия функций, индекс $M$ не вычисляется. Число таких случаев сообщается отдельно.

Индекс имеет вид
\begin{formulakey}
\begin{align}
M={}&w_{\mathrm{cc}}(1-\widehat{\mathrm{CC}})
+w_{\mathrm{dup}}(1-\widehat{\mathrm{Dup}})\notag\\
&+w_{\mathrm{warn}}(1-\widehat{\mathrm{Warn}})
+w_{\mathrm{doc}}\widehat{\mathrm{Doc}}.
\label{eq:maintainability}
\end{align}
\begin{symbols}
\sym{M}{сводная оценка свойств программы от 0 до~1}
\sym{w_{\mathrm{cc}},\,w_{\mathrm{dup}},\,w_{\mathrm{warn}},\,w_{\mathrm{doc}}}{веса сложности, дублирования, предупреждений и~документирования}
\sym{1-\widehat{\mathrm{CC}},\,1-\widehat{\mathrm{Dup}},\,1-\widehat{\mathrm{Warn}}}{вклады трёх признаков, для которых меньшие значения предпочтительны}
\sym{\widehat{\mathrm{Doc}}}{доля документированных публичных функций}
\end{symbols}
\end{formulakey}
Веса неотрицательны и~в~сумме равны единице, поэтому $M\in[0,1]$. При прочих равных индекс не уменьшается при снижении сложности, доли дубликатов и~числа предупреждений. Увеличение доли документированных функций также не уменьшает индекс. Таблица~\ref{tab:weights} содержит веса предложенной шкалы.

\begin{table}[!ht]
\caption{Предлагаемые начальные веса индекса свойств кода}
\label{tab:weights}
\centering\small\renewcommand{\arraystretch}{1.2}
\begin{tabularx}{\textwidth}{@{}lc>{\raggedright\arraybackslash}X@{}}
\toprule
Вес & Значение & Признак \\
\midrule
$w_{\mathrm{cc}}$ & 0,35 & Цикломатическая сложность \\
$w_{\mathrm{dup}}$ & 0,25 & Дублирование кода \\
$w_{\mathrm{warn}}$ & 0,20 & Предупреждения анализатора \\
$w_{\mathrm{doc}}$ & 0,20 & Наличие документации \\
\botrule
\end{tabularx}
\end{table}
\FloatBarrier

Веса выбраны для проверки чувствительности результатов. Стандарт ISO/IEC~25010:2023~\cite{iso25010} этих чисел не устанавливает. Оценки специалистов и~время сопровождения в~эксперименте не измерялись.

\begin{formulakey}
Для программы с~доступной оценкой свойств кода дополнительный балл равен
\[
S=qM^{\lambda},
\]
\begin{symbols}
\sym{S,\,q}{балл программы и~исход независимых тестов. При неуспехе $q=0$ и~балл равен нулю}
\sym{M}{сводная оценка свойств программы от 0 до~1}
\sym{\lambda}{вес этой оценки в~показателе, от 0 до~1. Исходное значение равно $0{,}5$}
\end{symbols}
\end{formulakey}
При $\lambda=0$ множитель $M^{\lambda}$ принимается равным единице, в~том числе для $M=0$.

Балл $S$ публикуется вместе с~основной долей успехов $Q$, всеми исходными компонентами и~долей кандидатов, для которых его удалось вычислить. В~пяти сериях, представленных далее, показатели $M$, $S$ и~VCR не измерялись. 

\subsection{Показатели обнаружения ошибок}
Полнота предварительной проверки $\mathrm{SCR}$ показывает, какую долю ошибочных кандидатов она выявила. Для её расчёта число отклонённых ошибочных кандидатов делится на число всех ошибочных кандидатов по независимой проверке. Вместе с~полнотой измеряется доля правильных кандидатов, которые предварительная процедура отклонила ошибочно. Эти два показателя позволяют оценить и~обнаружение ошибок, и~потерю правильных решений.

Показатель $\mathrm{RFR}$ равен доле применённых изменений, которые не были отменены за заранее установленный период. Для его расчёта нужны наблюдения за использованием программ после внесения изменений. Один запуск HumanEval таких данных не даёт.

Если в~выборке нет ошибочных кандидатов или применённых изменений, соответствующий знаменатель равен нулю. В~таком случае доля не определена. Это состояние сохраняется в~отчёте отдельно от нулевого значения измеренного показателя.


Показатели SCR и~RFR также не измерялись в~проведённых сериях.
